\section{Introduction}
\label{sec:intro}

A language model may need mathematical reasoning, code generation, scientific knowledge, and instruction following within the same deployment. Multi-teacher on-policy distillation (MOPD) transfers such capabilities by having specialized teachers score responses generated by a shared student \citep{ma2026mopdmultiteacheronpolicydistillation}. Because the student supplies the responses, its current behavior determines both the examples on which teachers provide feedback and the lengths over which that feedback is averaged.

Equal numbers of prompts from different domains need not give those domains equal weight in the loss. Open-MOPD identifies response-length imbalance as a source of unequal supervision \citep{gao2026openmopddiagnosingfixingcapability}. This raises a more specific question about the objective: does an averaging rule balance domains, responses within each domain, or both? These choices can be separated algebraically and compared empirically through the capabilities of the resulting student.

The parameter changes induced by teacher feedback raise a second question. Prior work finds concentrated changes in reinforcement learning \citep{mukherjee2025subnetworks} and distillation \citep{yu2026denseupdates,shen2026opdgeometry}. In a shared student, each teacher acts on parameters already influenced by other teachers. Concentration alone does not specify whether teachers affect the same locations. The relevant comparison therefore begins with a common student and optimizer state, and relates teacher-specific parameter changes to differences between teacher predictions.

A third question concerns the breadth of teacher feedback. A sampled-token loss and a loss over several vocabulary alternatives supervise the same generated prefix differently. Their vocabulary coverage does not directly determine how many model parameters change: the softmax couples logits, and the optimizer combines the current gradient with its history. Comparing both update concentration and capability connects the local learning signal to training outcomes.

We organize the paper around these three questions:
\begin{enumerate}
\item \textbf{Token balancing (Section~\ref{sec:normalization}).}
We derive how global-token, domain-token, and domain-response averaging weight domains and response lengths. The experiments show how these choices redistribute capability gains across domains.
\item \textbf{Update sparsity and teacher overlap (Section~\ref{sec:subnetworks}).}
We measure concentration of joint parameter changes, distinguishing individual steps from cumulative changes and optimizer weights from stored model weights. We formulate the teacher comparison at a shared student state and examine teacher-conditioned gradients in the appendix.
\item \textbf{Supervision density (Section~\ref{sec:density}).}
We compare sampled-token and top-$k$ supervision in single-teacher and joint training, and evaluate broader supervision at a common state. Local update concentration and online capability provide complementary measurements of these losses.
\end{enumerate}
